\documentclass[12pt, a4paper]{article}

% ── Shared preamble (src/lecturenotes.sty) ────────────────────────────────────
%   Packages, colors, boxes, \sol/\soleq and math macros all live in
%   src/*.sty. \input@path points LaTeX two folders up to find them.
\makeatletter
\def\input@path{{../../}}
\makeatother
\usepackage{lecturenotes}

% ── Document ──────────────────────────────────────────────────────────────────
\begin{document}

\begin{center}
  {\LARGE \textbf{Section 3.3: The Forward Pass}} \\[0.3em]
  {\large\itshape Computing a Prediction Step by Step} \\[0.5em]
  {\large AI-IM Class --- Lecture Notes}%
  \ifstudentmode
    \quad {\normalsize\color{gray}[Student Copy]}%
  \fi
\end{center}

\vspace{1em}
\hrule
\vspace{1.5em}

% ══════════════════════════════════════════════════════════════════════════════
The goal right now is to perform gradient descent on a simple feedforward network. The first step is to compute the \textbf{forward pass}, which means computing the network's prediction $\yhat$ for a given input $\bx$.
\section{The Network Architecture}
We will work with a \textbf{two-layer feedforward network} (one hidden layer)
designed to solve XOR:

\begin{figure}[H]
  \centering
  \begin{tikzpicture}[
      neuron/.style={circle, draw=mainblue, fill=boxbg,
          minimum size=1.1cm, thick, font=\small},
      input/.style={circle, draw=gray!60, fill=gray!15,
          minimum size=1.1cm, thick, font=\small},
      >=Stealth, node distance=0pt
    ]
    % ── Nodes ──────────────────────────────────────────────────────────────────
    \node[input]   (x1) at (0,  1.4) {$x_1$};
    \node[input]   (x2) at (0, -1.4) {$x_2$};

    \node[neuron]  (h1) at (4,  1.4) {$h_1$};
    \node[neuron]  (h2) at (4, -1.4) {$h_2$};

    \node[neuron]  (yh) at (8,  0)   {$\yhat$};

    % ── Connections: input → hidden ────────────────────────────────────────────
    \draw[->, thick] (x1) -- node[above, font=\scriptsize, pos=0.5]
    {$w^{(1)}_{11}=\phantom{-}0.5$} (h1);
    \draw[->, thick] (x1) -- node[above left, font=\scriptsize, pos=0.35]
    {$w^{(1)}_{21}=\phantom{-}0.5$} (h2);
    \draw[->, thick] (x2) -- node[below left, font=\scriptsize, pos=0.35]
    {$w^{(1)}_{12}=-0.5$} (h1);
    \draw[->, thick] (x2) -- node[below, font=\scriptsize, pos=0.5]
    {$w^{(1)}_{22}=\phantom{-}0.5$} (h2);

    % ── Connections: hidden → output ───────────────────────────────────────────
    \draw[->, thick] (h1) -- node[above, font=\scriptsize, pos=0.5]
    {$w^{(2)}_1=0.5$} (yh);
    \draw[->, thick] (h2) -- node[below, font=\scriptsize, pos=0.5]
    {$w^{(2)}_2=0.5$} (yh);

    % ── Bias annotations ───────────────────────────────────────────────────────
    \node[font=\scriptsize, gray!80!black, above=0.25cm of h1]
    {$b^{(1)}_1 = 0.5$};
    \node[font=\scriptsize, gray!80!black, below=0.25cm of h2]
    {$b^{(1)}_2 = 0.0$};
    \node[font=\scriptsize, gray!80!black, above=0.25cm of yh]
    {$b^{(2)} = 0.0$};

    % ── Layer labels ───────────────────────────────────────────────────────────
    \node[font=\small\itshape, gray] at (0,  3.1) {Input layer};
    \node[font=\small\itshape, gray] at (4,  3.1) {Hidden layer ($\sigmoid$)};
    \node[font=\small\itshape, gray] at (8,  3.1) {Output ($\sigmoid$)};
  \end{tikzpicture}
  \caption{A two-layer feedforward network with one hidden layer.}
  \label{fig:architecture}
\end{figure}

\vspace{0.5em}

The network has the following parameters:
\[
  \bW^{(1)} = \begin{bmatrix} 0.5 & -0.5 \\ 0.5 & \phantom{-}0.5 \end{bmatrix},
  \quad
  \bb^{(1)} = \begin{bmatrix} 0.5 \\ 0.0 \end{bmatrix},
  \quad
  \bW^{(2)} = \begin{bmatrix} 0.5 & 0.5 \end{bmatrix},
  \quad
  b^{(2)} = 0.0.
\]

Row $i$ of $\bW^{(1)}$ holds the weights \emph{feeding into} hidden unit $i$.
These weights are \textbf{not yet trained} --- we will start gradient descent
in Section 3.4.

% ══════════════════════════════════════════════════════════════════════════════
\section{Sigmoid Function}
In Section 3.1, we used the \textbf{step function} to compute the hidden
activations $h$. The step function is rarely used in practice because it is
flat everywhere except at the jump, so its derivative is $0$ (and undefined
at the jump itself).
Gradient descent (Section 3.2) moves the weights in the direction given by
the gradient. A zero gradient means a small change in a weight (almost) never changes
the output, so the network gets no signal telling it which way to move.

We therefore replace it with a smooth ``S-shaped'' version of the step,
the \textbf{sigmoid} function. Each neuron in the hidden and output layers uses
the sigmoid activation:
\[
  \sigmoid(z) = \frac{1}{1 + e^{-z}}, \qquad \sigmoid(z) \in (0, 1).
\]

The computation inside every neuron follows two sub-steps:
\begin{enumerate}[nosep]
  \item \textbf{Pre-activation} $z$: a weighted sum of inputs plus bias.
  \item \textbf{Activation} $h = \sigmoid(z)$: squeeze $z$ into $(0,1)$.
\end{enumerate}

Making both sub-steps explicit, every neuron in our network splits into a
$z$ node followed by an activation node:

\begin{figure}[H]
  \centering
  \begin{tikzpicture}[
      neuron/.style={circle, draw=mainblue, fill=boxbg,
          minimum size=1.1cm, thick, font=\small},
      preact/.style={circle, draw=mainblue, fill=white,
          minimum size=1.1cm, thick, font=\small},
      input/.style={circle, draw=gray!60, fill=gray!15,
          minimum size=1.1cm, thick, font=\small},
      unitbox/.style={draw=gray!60, dashed, rounded corners=6pt},
      >=Stealth, node distance=0pt
    ]
    % ── Nodes ──────────────────────────────────────────────────────────────────
    \node[input]   (x1) at (0,    1.4) {$x_1$};
    \node[input]   (x2) at (0,   -1.4) {$x_2$};

    \node[preact]  (z1) at (3.6,  1.4) {$z^{(1)}_1$};
    \node[neuron]  (h1) at (5.6,  1.4) {$h_1$};
    \node[preact]  (z2) at (3.6, -1.4) {$z^{(1)}_2$};
    \node[neuron]  (h2) at (5.6, -1.4) {$h_2$};

    \node[preact]  (zo) at (9.2,  0)   {$z^{(2)}$};
    \node[neuron]  (yh) at (11.2, 0)   {$\yhat$};

    % ── Dashed boxes: one box = one neuron ─────────────────────────────────────
    \draw[unitbox] (2.85,  0.65) rectangle (6.35,  2.15);
    \draw[unitbox] (2.85, -2.15) rectangle (6.35, -0.65);
    \draw[unitbox] (8.45, -0.75) rectangle (11.95, 0.75);

    % ── Connections: input → hidden pre-activations ────────────────────────────
    \draw[->, thick] (x1) -- node[above, font=\scriptsize, pos=0.5]
    {$w^{(1)}_{11}$} (z1);
    \draw[->, thick] (x1) -- node[sloped, above, font=\scriptsize, pos=0.28]
    {$w^{(1)}_{21}$} (z2);
    \draw[->, thick] (x2) -- node[sloped, below, font=\scriptsize, pos=0.28]
    {$w^{(1)}_{12}$} (z1);
    \draw[->, thick] (x2) -- node[below, font=\scriptsize, pos=0.5]
    {$w^{(1)}_{22}$} (z2);

    % ── Activation arrows z → h ────────────────────────────────────────────────
    \draw[->, thick, mainblue] (z1) -- node[above, font=\small] {$\sigmoid$} (h1);
    \draw[->, thick, mainblue] (z2) -- node[above, font=\small] {$\sigmoid$} (h2);
    \draw[->, thick, mainblue] (zo) -- node[above, font=\small] {$\sigmoid$} (yh);

    % ── Connections: hidden → output pre-activation ────────────────────────────
    \draw[->, thick] (h1) -- node[above right, font=\scriptsize, pos=0.45]
    {$w^{(2)}_1$} (zo);
    \draw[->, thick] (h2) -- node[below right, font=\scriptsize, pos=0.45]
    {$w^{(2)}_2$} (zo);

    % ── Bias annotations ───────────────────────────────────────────────────────
    \node[font=\scriptsize, gray!80!black] at (3.6,  2.45) {$+\,b^{(1)}_1$};
    \node[font=\scriptsize, gray!80!black] at (3.6, -2.45) {$+\,b^{(1)}_2$};
    \node[font=\scriptsize, gray!80!black] at (9.2,  1.05) {$+\,b^{(2)}$};

    % ── Layer labels ───────────────────────────────────────────────────────────
    \node[font=\small\itshape, gray] at (0,    3.1) {Input layer};
    \node[font=\small\itshape, gray] at (4.6,  3.1) {Hidden layer};
    \node[font=\small\itshape, gray] at (10.2, 3.1) {Output layer};
  \end{tikzpicture}
  \caption{The same network with each neuron (dashed box) split into its two
    sub-steps. The white node holds the pre-activation $z$, the weighted sum
    of the incoming values plus the bias. The blue node holds the activation
    obtained by applying the sigmoid, $h = \sigmoid(z)$ (or $\yhat = \sigmoid(z^{(2)})$
    at the output).}
  \label{fig:forward-pass-z}
\end{figure}

\begin{notebox}
  \textbf{Computing $\sigmoid$ with NumPy.}
  You can evaluate the sigmoid in Python with:
\end{notebox}
\begin{codebox}
import numpy as np

def sigmoid(x):
    return 1 / (1 + np.exp(-x))
\end{codebox}

% ══════════════════════════════════════════════════════════════════════════════
\section{The Forward Pass}

We trace the network with the XOR input $\bx = (x_1, x_2) = (0, 1)$,
whose target output is $y = 1$.

% ──────────────────────────────────────────────────────────────────────────────
\begin{stepbox}{Step 1 \quad Input $\to$ Hidden pre-activations}

  Compute the weighted sum entering each hidden neuron:

  \[
    z^{(1)}_1
    = w^{(1)}_{11}\,x_1 + w^{(1)}_{12}\,x_2 + b^{(1)}_1
    = (0.5)(0) + (-0.5)(1) + 0.5
    = \sol[3cm]{0.0}
  \]

  \[
    z^{(1)}_2
    = w^{(1)}_{21}\,x_1 + w^{(1)}_{22}\,x_2 + b^{(1)}_2
    = (0.5)(0) + (0.5)(1) + 0.0
    = \sol[3cm]{0.5}
  \]
\end{stepbox}

\vspace{0.8em}

% ──────────────────────────────────────────────────────────────────────────────
\begin{stepbox}{Step 2 \quad Hidden activations}

  Apply the sigmoid to each pre-activation (use the \texttt{sigmoid} function above):

  \[
    h_1 = \sigmoid\!\left(z^{(1)}_1\right) = \sigmoid(0.0) = \sol[3cm]{0.500}
  \]

  \[
    h_2 = \sigmoid\!\left(z^{(1)}_2\right) = \sigmoid(0.5) = \sol[3cm]{0.622}
  \]
\end{stepbox}

\vspace{0.8em}

% ──────────────────────────────────────────────────────────────────────────────
\begin{stepbox}{Step 3 \quad Hidden $\to$ Output pre-activation}

  The output neuron collects the two hidden activations:

  \[
    z^{(2)}
    = w^{(2)}_1\,h_1 + w^{(2)}_2\,h_2 + b^{(2)}
    = (0.5)(0.500) + (0.5)(0.622) + 0.0
    = 0.250 + 0.311
    = \sol[3cm]{0.561}
  \]
\end{stepbox}

\vspace{0.8em}

% ──────────────────────────────────────────────────────────────────────────────
\begin{stepbox}{Step 4 \quad Output activation}

  \[
    \yhat = \sigmoid\!\left(z^{(2)}\right) = \sigmoid(0.561)
  \]

  Using the \texttt{sigmoid} function above, \texttt{sigmoid(0.561)} gives:

  \[
    \yhat \approx \sol[3cm]{0.637}
  \]
\end{stepbox}

% ══════════════════════════════════════════════════════════════════════════════
\section{Computing the Loss}

The target for input $(0,1)$ is $y = 1$.
Using Mean Squared Error:

\[
  J = (y - \yhat)^2 = (1 - 0.637)^2 = (0.363)^2 \approx \sol[3cm]{0.132}
\]

The network is wrong: it outputs $0.637$ when it should output $1$.
The loss $J = 0.132$ tells us how far off we are.

\begin{notebox}
  A perfect prediction ($\yhat = 1$) would give $J = 0$.
  A blind guess of $\yhat = 0.5$ would give $J = (0.5)^2 = 0.25$.
  Our untrained network ($J = 0.132$) happens to be a bit better than random
  --- but only because we are looking at a single input.
\end{notebox}

% ══════════════════════════════════════════════════════════════════════════════
\section{Summary: What the Forward Pass Gives Us}
The network has the following parameters:
\[
  \bW^{(1)} = \begin{bmatrix} 0.5 & -0.5 \\ 0.5 & \phantom{-}0.5 \end{bmatrix},
  \quad
  \bb^{(1)} = \begin{bmatrix} 0.5 \\ 0.0 \end{bmatrix},
  \quad
  \bW^{(2)} = \begin{bmatrix} 0.5 & 0.5 \end{bmatrix},
  \quad
  b^{(2)} = 0.0.
\]
\begin{figure}[H]
  \centering
  \begin{tikzpicture}[
      neuron/.style={circle, draw=mainblue, fill=boxbg,
          minimum size=1.1cm, thick, font=\small},
      preact/.style={circle, draw=mainblue, fill=white,
          minimum size=1.1cm, thick, font=\small},
      input/.style={circle, draw=gray!60, fill=gray!15,
          minimum size=1.1cm, thick, font=\small},
      unitbox/.style={draw=gray!60, dashed, rounded corners=6pt},
      >=Stealth, node distance=0pt
    ]
    % ── Nodes ──────────────────────────────────────────────────────────────────
    \node[input]   (x1) at (0,    1.4) {$x_1$};
    \node[input]   (x2) at (0,   -1.4) {$x_2$};

    \node[preact]  (z1) at (3.6,  1.4) {$z^{(1)}_1$};
    \node[neuron]  (h1) at (5.6,  1.4) {$h_1$};
    \node[preact]  (z2) at (3.6, -1.4) {$z^{(1)}_2$};
    \node[neuron]  (h2) at (5.6, -1.4) {$h_2$};

    \node[preact]  (zo) at (9.2,  0)   {$z^{(2)}$};
    \node[neuron]  (yh) at (11.2, 0)   {$\yhat$};

    % ── Dashed boxes: one box = one neuron ─────────────────────────────────────
    \draw[unitbox] (2.85,  0.65) rectangle (6.35,  2.15);
    \draw[unitbox] (2.85, -2.15) rectangle (6.35, -0.65);
    \draw[unitbox] (8.45, -0.75) rectangle (11.95, 0.75);

    % ── Connections: input → hidden pre-activations ────────────────────────────
    \draw[->, thick] (x1) -- node[above, font=\scriptsize, pos=0.5]
    {$w^{(1)}_{11}$} (z1);
    \draw[->, thick] (x1) -- node[sloped, above, font=\scriptsize, pos=0.28]
    {$w^{(1)}_{21}$} (z2);
    \draw[->, thick] (x2) -- node[sloped, below, font=\scriptsize, pos=0.28]
    {$w^{(1)}_{12}$} (z1);
    \draw[->, thick] (x2) -- node[below, font=\scriptsize, pos=0.5]
    {$w^{(1)}_{22}$} (z2);

    % ── Activation arrows z → h ────────────────────────────────────────────────
    \draw[->, thick, mainblue] (z1) -- node[above, font=\small] {$\sigmoid$} (h1);
    \draw[->, thick, mainblue] (z2) -- node[above, font=\small] {$\sigmoid$} (h2);
    \draw[->, thick, mainblue] (zo) -- node[above, font=\small] {$\sigmoid$} (yh);

    % ── Connections: hidden → output pre-activation ────────────────────────────
    \draw[->, thick] (h1) -- node[above right, font=\scriptsize, pos=0.45]
    {$w^{(2)}_1$} (zo);
    \draw[->, thick] (h2) -- node[below right, font=\scriptsize, pos=0.45]
    {$w^{(2)}_2$} (zo);

    % ── Bias annotations ───────────────────────────────────────────────────────
    \node[font=\scriptsize, gray!80!black] at (3.6,  2.45) {$+\,b^{(1)}_1$};
    \node[font=\scriptsize, gray!80!black] at (3.6, -2.45) {$+\,b^{(1)}_2$};
    \node[font=\scriptsize, gray!80!black] at (9.2,  1.05) {$+\,b^{(2)}$};

    % ── Layer labels ───────────────────────────────────────────────────────────
    \node[font=\small\itshape, gray] at (0,    3.1) {Input layer};
    \node[font=\small\itshape, gray] at (4.6,  3.1) {Hidden layer};
    \node[font=\small\itshape, gray] at (10.2, 3.1) {Output layer};
    \node[font=\small\itshape, gray] at (13.6, 3.1) {Loss};

    % ── Loss node and target ───────────────────────────────────────────────────
    \node[circle, draw=red!70!black, fill=warnbg, minimum size=1.1cm,
      thick, font=\small] (J) at (13.6, 0) {$J$};
    \node[input] (y) at (13.6, -2.0) {$y$};
    \draw[->, thick, red!70!black] (yh) -- (J);
    \draw[->, thick, red!70!black] (y)  -- (J);
    \node[font=\scriptsize, red!70!black] at (13.6, 0.95) {$J = (y - \yhat)^2$};
    \node[font=\scriptsize, gray!80!black, left=0.1cm of y] {target};
  \end{tikzpicture}
  \caption{The complete forward pass, from the inputs $x_1, x_2$ through the
    pre-activations $z$ and activations $h$ to the prediction $\yhat$. The
    prediction is then compared with the target $y$ to give the loss
    $J = (y - \yhat)^2$. Backpropagation (Section 3.4) traverses this same
    graph in reverse, from $J$ back to every weight.}
  \label{fig:forward-pass-z-with-loss}
\end{figure}

\begin{center}
  \renewcommand{\arraystretch}{1.4}
  \begin{tabular}{lll}
    \toprule
    Quantity                & Symbol         & Value      \\
    \midrule
    Input                   & $(x_1,\, x_2)$ & $(0,\; 1)$ \\
    Target                  & $y$            & $1$        \\
    \midrule
    Hidden pre-activation 1 & $z^{(1)}_1$    & $0.0$      \\
    Hidden pre-activation 2 & $z^{(1)}_2$    & $0.5$      \\
    Hidden activation 1     & $h_1$          & $0.500$    \\
    Hidden activation 2     & $h_2$          & $0.622$    \\
    \midrule
    Output pre-activation   & $z^{(2)}$      & $0.561$    \\
    Prediction              & $\yhat$        & $0.637$    \\
    \midrule
    Loss                    & $J$            & $0.132$    \\
    \bottomrule
  \end{tabular}
\end{center}

\vspace{0.5em}

\begin{notebox}
  \textbf{These numbers are the starting point for Section 3.4.}
  Backpropagation will work \emph{backwards} through this same network to
  compute $\partial J / \partial w$ for every weight --- which is exactly what
  gradient descent needs to take an update step.
\end{notebox}

\end{document}
